\section{Introduction}

Instrumental variables (IVs) are a standard tool for identifying causal effects when treatment assignment is endogenous in observational studies. Under the classical IV assumptions, \citet{angrist1995identification} show that the Local Average Treatment Effect (LATE), defined as the average effect for the subpopulation of compliers, is point-identified. However, many policy questions are not about compliers themselves; they ask how outcomes would change under a counterfactual policy that shifts treatment take-up in the broader population. Such questions are naturally captured by Policy-Relevant Treatment Effects (PRTEs), which are central targets in the generalized Roy and Marginal Treatment Effect (MTE) framework \citep{heckman1999local,heckman2005structural,heckman2006understanding}. This framework encompasses a wide range of causal estimands, including the Average Treatment Effect (ATE); empirical policy evaluations within it include \citet{carneiro2010evaluating,carneiro2011estimating}. Unlike LATE, PRTEs are generally not point-identified when the instrument induces only limited support in the treatment propensity.

%Classical applications include the causal return to education \citep{angrist1991does,card1993using}, the earnings effects of military service \citep{angrist1990lifetime}, the long-run impact of institutions on economic development \citep{acemoglu2001colonial}, and policy interventions with imperfect treatment compliance \citep{angrist1995identification}. 

This limited-support problem is central in the generalized Roy model \citep{heckman2005structural}, where treatment selection follows a threshold-crossing mechanism formally equivalent to the classical monotonicity condition \citep{angrist1995identification,vytlacil2002independence}. In this framework, PRTEs can be written as weighted averages of the MTE, so identification depends on how much of the latent resistance distribution is reached by the propensity score induced by the instrument. When the propensity score lacks full support on the unit interval, PRTEs and related quantities such as the ATE remain only partially identified, and a long line of work has studied such identified sets under various assumptions \citep{manski1990nonparametric,manski1997monotone,manski_pepper2000monotone,heckman2001instrumental,magne2018using}. 

Despite its importance, deriving tight, closed-form bounds for general PRTEs remains a significant challenge. Existing strategies fall into three broad approaches: moment-relaxation over MTR functions \citep{magne2018using,mogstad2018identification}, second-order stochastic dominance constraints on the MTE function solved by rearrangement \citep{marx2024sharp}, and infinite-dimensional linear programming over the joint law of outcome and latent resistance solved by sieve approximation \citep{han_yang2024computational}. Each faces limitations---discarded distributional information, restricted scope to weights on the complier interval under a binary instrument, lack of closed-form solutions or incompatibility with continuous instruments---that we discuss in detail in \cref{subsec:related_work}.

In this paper, we develop a unified Optimal Transport (OT) framework for partial identification in the generalized Roy model and use it to derive sharp, closed-form bounds for PRTEs. Specifically, we recast PRTE partial identification as a Constrained Conditional Optimal Transport (CCOT) problem, a structured instance within the broader OT framework that is equivalent to the infinite-dimensional LP of \citet{han_yang2024computational}. Rather than summarizing the observed distributions through moment conditions or imposing distributional constraints on conditional mean (MTE) functions, we place the distributional constraints directly on the joint conditional law of the potential outcome and latent resistance. Then, we exploit the OT structure to prove that this heavily constrained, multidimensional CCOT problem analytically reduces to separable one-dimensional OT problems. This reduction completely bypasses computationally intensive linear programming and sieve approximation, yielding explicit, closed-form solutions for the sharp bounds. 

Our main contributions can be summarized as follows.

\begin{enumerate}
    \item \textbf{A Unified OT Framework for the Generalized Roy Model:} We formulate partial identification in the generalized Roy model as a CCOT problem. The framework imposes distributional constraints directly at the joint-law level, and we prove that the constraints decompose analytically into separable marginal constraints.
    \item \textbf{Analytic Decomposition and Closed-Form PRTE Bounds:} Leveraging the decomposition, we prove the CCOT problem decomposes into separable one-dimensional OT problems with product costs, yielding explicit, closed-form sharp bounds and bypassing the sieve approximation required by existing computational approaches.
    \item \textbf{Semiparametric Estimation and Inference:} We develop estimation results for the binary-treatment bounds. For discrete instruments, we leverage Double Machine Learning (DML) to construct a debiased estimator, achieving $\sqrt{n}$-consistency and asymptotic normality. For continuous instruments, we explicitly characterize the nonparametric convergence rates.
    \item \textbf{Empirical Validation:} We validate our theoretical results through both synthetic simulations and a real-world empirical application, demonstrating how our closed-form bounds tighten the identified set relative to the moment-relaxation methods.
\end{enumerate}

Replacing threshold crossing by a conditional quantile representation of the selection equation yields sharp closed-form bounds for continuous and multi-valued treatments, under a rank-overlap condition and an algebraic condition on the instrument-induced latent intervals, respectively. A companion note develops those results in full; the present paper treats the binary case throughout.



\paragraph{Organization of this paper.} \cref{sec:pre} introduces the setup and optimal transport background; \cref{sec:pid} presents the main identification theory and closed-form bounds; \cref{sec:est_inf} develops estimation and inference; and \cref{sec:simulation} reports simulation and empirical results. The main identification and inference proofs, the algorithms, and the orthogonal score formulas are in \supplocation\suppnumnote; \oalocation\ collects the notation table, worked examples, supporting derivations, and full simulation and empirical documentation.

\subsection{Related Work}\label{subsec:related_work}
\paragraph{Local IV and MTE.}
The Marginal Treatment Effect (MTE) of \citet{heckman1999local,heckman2005structural} is the building block of our identification analysis: it is the treatment effect at a given quantile of unobserved resistance, and all standard IV estimands---LATE, ATE, ATT, and PRTE---are weighted averages of it \citep{heckman2005structural, heckman2006understanding}. \citet{vytlacil2002independence} establishes the equivalence between the threshold-crossing selection model and the monotonicity assumption of \citet{angrist1995identification}. Point identification proceeds via the Local Instrumental Variables (LIV) estimand \citep{heckman1999local}, which requires full propensity score support over $[0,1]$---the condition that fails under limited instrument variation, the regime our paper addresses.

% \paragraph{IV with General Treatments.} \citet{angrist1995two} extends the LATE framework to ordered discrete treatments and \citet{heckman2007econometric} extends the MTE framework to an ordered-choice model; \citet{kirkeboen2016field} treats multiple unordered discrete choices by accounting for self-selection on unobserved comparative advantage, and \citet{lee2018identifying} develop identification results for multivalued treatments. For continuous treatments, \citet{florens2008identification} and \citet{imbens2009identification} use control-function methods to identify, respectively, the ATE and ATT and structural equations with continuous endogenous variables. These results deliver point identification of specific causal quantities; we instead target partial identification of PRTEs, and a companion note extends our results to the continuous and multi-valued settings.

\paragraph{Partial Identification.} \citet{manski1990nonparametric,manski1997monotone,manski_pepper2000monotone} pioneered partial identification, providing closed-form bounds for the IV problem under a range of assumptions; as a special case, our bound recovers the classical Manski bound for the ATE. In the econometric literature, partial identification is often cast as estimating the identified set defined by moment inequalities \citep{chernozhukov_hong_tamer2007}; see \citet{tamer2010partial} for a comprehensive survey.

\citet{balke_pearl1997bounds} develop the linear programming approach for partial identification, which has recently been revisited and generalized using modern techniques \citep{balazadeh_meresht_syrgkanis_krishnan2022,guo_yin_wang_jordan2022,duarte2024automated,levis2025covariate}; in particular, \citet{levis2025covariate} use covariates to tighten the classical bound of \citet{balke_pearl1997bounds} and develop estimation theory in the IV setting. Several works connect partial identification to robust optimization \citep{guo_yin_wang_jordan2022,balazadeh_meresht_syrgkanis_krishnan2022,fan_pass_shi2025,tan2024consistency}, and, in settings without unmeasured confounding, the problem can be formulated as an optimal transport problem with marginal constraints derived from the observed-data distribution \citep{ji_lei_spector2024,lin2025estimation,lin2025tightening}; see \citet{gunsilius2025primer} for a recent primer. In related but distinct settings, \citet{fan_guerre_zhu2017partial} study partial identification of functionals of the joint distribution of potential outcomes when the conditional marginals are identified, and \citet{oberreynolds2023estimating} characterizes sharp bounds on functionals of the joint complier outcome distribution via optimal transport under the binary-instrument LATE framework of \citet{angrist1995identification}. The first bounds functionals of the joint law from identified marginals, the second an identified subpopulation; neither addresses the PRTE weights we study here.

\paragraph{PRTE Partial Identification.} Three strategies have emerged for PRTE partial identification under limited support. \citet{magne2018using} optimize over MTR functions subject to moment constraints \citep{mogstad2018identification}. This moment relaxation discards distributional information beyond the first moment. \citet{han_yang2024computational} preserve full distributional content by formulating an infinite-dimensional LP over the joint law of latent states, solved by sieve approximation for discrete instruments. \citet{marx2024sharp} generalizes the approach of \citet{magne2018using} by imposing full distributional content on the MTE function via a second-order stochastic dominance (SSD) constraint, obtaining sharp closed-form bounds by rearrangement; the paper's explicit theorems cover weights on the complier interval under a binary instrument, with extensions beyond compliers requiring rank similarity.

Our work differs from all three in scope: we impose distributional constraints directly on the joint conditional law of the potential outcome and latent resistance rather than on a conditional mean, cast the resulting program as a CCOT problem equivalent to the LP of \citet{han_yang2024computational}, and show it reduces analytically to separable one-dimensional OT problems with product costs. This delivers (i) sharp closed-form bounds for the full PRTE on all of $[0,1]$ under multi-valued instruments via an interval decomposition across $K\!+\!1$ compliance types, including boundary intervals where only one potential outcome marginal is identified, (ii) extensions to continuous and mixed instruments with sub-problems indexed by the propensity score continuum, and (iii) discrete-instrument inference together with continuous-instrument rate results, which are limited in either \citet{marx2024sharp} or \citet{han_yang2024computational}. On the complier interval under a binary instrument, the CCOT solution and the SSD rearrangement of \citet{marx2024sharp} coincide. Beyond PRTE, our sharpness result characterizes the identified set at the level of the full joint latent law and sheds light on the joint-identification question raised in \citet{marx2024sharp} through its reduction to independent one-dimensional marginal constraints.

\noindent\textit{Remark (notation).} $K$ is a population count in propensity decompositions and a chosen function when written as the local-linear smoothing kernel $K(t)$. In Theorem~\ref{thm:discrete_lower_bound}: $K=|\{p(z):z\in\mathcal Z,\ p(z)>0\}|$, excluding the structural boundary $p_0=0$. In Theorem~\ref{thm:continuous_lower_bound}: $K$ is the number of ordered maximal closed intervals in $p(\mathcal Z)$ after merging overlaps, including degenerate singleton intervals and any interval containing zero. In Proposition~\ref{prop:discrete_reduction}, Proposition~\ref{prop:J_specialization}, Lemma~\ref{lem:orthogonality_discrete}, Theorem~\ref{thm:dml_normality}: $K$ counts the distinct positive baseline propensity levels at each $x$ and does not depend on $x$; $k=0,\ldots,K-1$ indexes the identified intervals and $k=K$ indexes the residual tail. In Corollary~\ref{cor:continuous_primitive_rate}: $K(t)$ is the base one-dimensional smoothing kernel in $K_{h_n}(t)=K(t/h_n)/h_n$, rather than a propensity-level count.


% Crucially, while \citet{marx2024sharp} imposes distributional constraints only on the conditional mean (the MTE function) and does not characterize the identified set at the level of the joint latent distribution, our sharpness result (\cref{prop:sharpness}) directly characterizes the identified set of the full joint latent law of $(Y(0), Y(1), U, X)$.
